% Generated draft: requires mathematical and source review.
\documentclass[11pt]{article}
\usepackage[margin=25mm]{geometry}
\usepackage{fontspec}
\setmainfont{Latin Modern Roman}
\usepackage{amsmath,amssymb,mathtools,unicode-math}
\setmathfont{Latin Modern Math}
\usepackage{enumitem,needspace,xurl,hyperref}
\setlength{\parindent}{0pt}
\setlength{\parskip}{5pt}
\setlength{\emergencystretch}{3em}
\newcommand{\N}{\mathbb{N}}
\newcommand{\Z}{\mathbb{Z}}
\newcommand{\Q}{\mathbb{Q}}
\newcommand{\R}{\mathbb{R}}
\newcommand{\C}{\mathbb{C}}
\newcommand{\F}{\mathbb{F}}
\newcommand{\A}{\mathbb{A}}
\newcommand{\PP}{\mathbb{P}}
\newcommand{\E}{\mathbb{E}}
\newcommand{\eps}{\varepsilon}
\newcommand{\dd}{\,\mathrm{d}}
\begin{document}
\section*{UnsolvedMath ID: 700685}
\textbf{Source problem code:} catalogue entry
\subsection{Definitions and mathematical statement}
For a family of groups \((G_i)_{i\in I}\), their free product \(G=\mathop{\ast}_{i\in I}G_i\) is a group equipped with homomorphisms \(\iota_i:G_i\to G\) such that, for every group \(H\) and family of homomorphisms \(f_i:G_i\to H\), there exists a unique homomorphism \(f:G\to H\) satisfying \(f\circ\iota_i=f_i\) for every \(i\in I\). This notation does not impose a scope on the source problem: the supplied record does not specify restrictions on the factors or the index set.

The requested task is to find a geometrical justification of the Whitehead method for free products, analogous to Whitehead's own geometrical substantiation of the method for free groups, and more visual than the treatment of D. J. Collins and H. Zieschang in the two cited 1984 papers.

The supplied record does not define the version of the Whitehead method intended here, its inputs, its admissible transformations, or its asserted conclusions. It also does not specify what constitutes a geometrical justification or a more visual treatment. Consequently, this is a qualitative construction and exposition problem; a fully quantified correctness statement requires identifying the intended method and its hypotheses from the cited literature.

D. J. Collins' supplied comment reports a partial solution in D. McCullough and A. Miller, \emph{Symmetric automorphisms of free products} (1996). The record does not describe the scope of that partial solution or identify the precise remaining task.
\subsection{Short English statement}
Find a geometric explanation of the Whitehead method for free products, comparable to Whitehead's explanation for free groups and more visual than the cited Collins--Zieschang treatments. The exact method and hypotheses require clarification. Collins reports a partial solution by McCullough and Miller, whose scope is not specified in the supplied record.
\subsection{Sources}
\begin{itemize}
\item Local UnsolvedMath registry, permanent ID \(700685\), exact source problem code \texttt{catalogue entry}, title \emph{Kourovka Notebook Problem 10.70}. The supplied registry record labels the problem open; this label is not an independently verified mathematical assessment.
\item \emph{The Kourovka Notebook (Unsolved Problems in Group Theory)}, original problem 10.70, proposed by H. Zieschang in the 10th issue (1986). The supplied background also identifies the 21st issue edition (2026), \texttt{arXiv:1401.0300v48}. The collection description names E. I. Khukhro and V. D. Mazurov as editors.
\item D. J. Collins and H. Zieschang, \emph{Math. Z.} \textbf{185}, no. 4 (1984), 487--504; and \textbf{186}, no. 3 (1984), 335--361. These are the treatments used as the source problem's comparison for visual clarity. Article titles are not supplied.
\item D. McCullough and A. Miller, \emph{Symmetric automorphisms of free products}, Mem. Amer. Math. Soc. \textbf{582} (1996), cited in D. J. Collins' comment as providing a partial solution. The extent of that solution is not supplied.
\item Whitehead's geometrical substantiation for free groups is referenced descriptively in the problem; no specific bibliographic citation for it is supplied.
\end{itemize}
\end{document}
